% Conversão fiel da participação como cliente Hoortech; origem imutável RELATORIO_ATUAL.odt.
\hypertarget{__RefHeading___Toc4470_31181766333}{}

% SOURCE text-h-0033
\refstepcounter{subsection}
\subsection*{6.2.4 Sprint 2}
\label{cliente-hoortech:text-h-0033}
\addcontentsline{toc}{subsection}{6.2.4 Sprint 2}

% SOURCE text-p-0728
A entrega da Sprint 2 foi extremamente animadora, refletindo um grande avanço no desenvolvimento do projeto. Na ocasião, os alunos finalizaram completamente o front-end, enquanto o back-end se encontrava praticamente concluído, faltando apenas a fase de treinamento do modelo de reconhecimento de sinais. Eles já haviam planejado dedicar as próximas duas sprints exclusivamente ao treinamento do modelo e à preparação do dataset.\par

% SOURCE text-p-0729
A equipe iniciou também a implementação de pequenas Provas de Conceito (POCs), com algumas já deployadas na AWS. Um dos principais desafios relatados pelos alunos foi o reconhecimento de sinais que envolvem movimento contínuo, um problema semelhante ao que enfrentamos na Hackatona de 2023. Dada essa limitação, acordamos que o modelo inicial cobriria o alfabeto completo, com exceção das letras que requerem movimento. Em seguida, o próximo passo seria o desenvolvimento do reconhecimento de números.\par

% SOURCE text-p-0730
Outro ponto importante foi a questão das requisições na AWS. Os alunos relataram dificuldades em obter a frequência ideal de requisições, pois o sistema permitia apenas uma requisição a cada 100 milissegundos. Realizar mais requisições em intervalos menores acarretaria em custos elevados, inviabilizando a operação dentro do escopo do projeto atual. Entendemos que essa limitação representa um obstáculo técnico significativo e não há uma solução imediata.\par

% SOURCE text-p-0731
Dado o alto grau de complexidade do projeto, eu e minha equipe oferecemos nosso apoio para ajudar com algumas questões relacionadas ao modelo. Apesar dos desafios, ficou evidente que os alunos estão bem alinhados e preparados para lidar com as próximas etapas. Com uma comunicação eficaz e trabalho colaborativo, ficamos confiantes de que esse projeto de alta complexidade estava progredindo de forma consistente e promissora.\par

